\section{Setting and Nearest Approaches}
\label{sec:prior}

\paragraph{Setting}
We use electric distribution as the running example. A \emph{device} in this paper is the
enforcement point that decides whether to execute a remote command, for example a module
in a protection relay, a remote terminal unit or a gateway. Where such a point would sit,
and which operations must continue when a control room is isolated, are open questions
for people who operate grids (\cref{sec:open}). We consulted no such person.

\paragraph{Sources and how we read them}
We read the relevant sections of five sources in full: URSA, Nerio, Spire, the report by
Liskov and Rodrigues, and the paper by Shakarami and Sandhu. Two recent articles, HRCred
and the work by Stantchev, refused automated download. For them we have only titles,
abstracts and metadata, and we claim nothing about their content.
\Cref{tab:positioning} summarizes. The positioning record in the artifact gives page
references for each statement below.

\paragraph{Threshold-issued expiring credentials}
In URSA~\cite{luo2004ursa}, a ticket is valid if it is certified and unexpired. Its
neighbours renew it collectively with a threshold signature, and a node that they see
misbehave is denied renewal. Removal therefore needs no message to the device. URSA has
no notion of a single holder of authority or of an exclusive successor. HRCred
(abstract only)~\cite{zhou2026hrcred} combines threshold issuance with revocation epochs
and lifetime bounds for device credentials. It could overlap with this seed, and reading
it is open work (\cref{sec:open}). The same is true of Stantchev~\cite{stantchev2026survives},
whose title announces authorization that survives disconnection.

\paragraph{Exclusive leases}
In Nerio~\cite{vanrenesse2011nerio}, each process keeps an assignee and a finish time, and
never changes its assignee before the finish time. A process is the leader if a quorum
grants it a lease, so two leaders cannot overlap. The fault model is crash failures, and
the leader itself checks that its lease is still valid before it sends an edict. Leases in
general trade the length of a term against the effect of a failure~\cite{gray1989leases}.
A hostile holder ignores its own expiry. Our comparator \texttt{xl} (\cref{sec:comparisons})
therefore adds signatures, checks at the device and Byzantine quorums to Nerio's grantor
rule. This adaptation is ours, and Nerio's authors did not evaluate it.

\paragraph{Intrusion-tolerant SCADA}
Spire replicates the SCADA master with $3f + 2k + 1$ replicas. It tolerates $f$ compromised
replicas, one disconnected site and one replica in recovery~\cite{babay2018spire}. Its
authors write that their ``threat model does not cover a compromised HMI or a rogue
operator''~\cite{babay2018spire}. Our question is about that case. An exclusion in one
paper does not show that no other work solves it.

\paragraph{Residual effect after removal}
Liskov and Rodrigues bound the effect of a Byzantine client after an administrator removes
it: it can leave at most a bounded number of ``lurking writes''~\cite{liskov2005tr}. Their
replicas enforce the bound and can be reached. In our setting the devices decide alone,
and the bound is a time after withdrawal.

\paragraph{Freshness at a decision point}
Shakarami and Sandhu show that stale attribute values at a decision point cause both
safety and availability violations, and that refreshing values helps
both~\cite{shakarami2019refresh}. It is the same tension, for a different object.

\begin{table}[t]
\caption{Nearest approaches, from the parts of each source we read. ``Exclusive'' means a
rule that prevents overlapping authority of two holders.}
\label{tab:positioning}
\centering
\footnotesize
\begin{tabularx}{\columnwidth}{@{}>{\raggedright\arraybackslash}p{0.2\columnwidth}>{\raggedright\arraybackslash}X>{\raggedright\arraybackslash}p{0.19\columnwidth}ll@{}}
\toprule
Approach & Faults & Authority ends by & Excl. & Read \\
\midrule
URSA~\cite{luo2004ursa} & Byzantine nodes below a threshold & expiry, no renewal & no & full \\
HRCred~\cite{zhou2026hrcred} & unknown & revocation epochs, lifetimes & ? & abstract \\
Nerio~\cite{vanrenesse2011nerio} & crash & lease expiry & yes & full \\
Spire~\cite{babay2018spire} & $f$ replicas, one site cut & (operator trusted) & --- & full \\
Liskov, Rodrigues~\cite{liskov2005tr} & Byzantine clients and replicas & administrator & --- & full \\
This seed & $f$ officers, hostile holder & expiry; epoch grants & yes & --- \\
\bottomrule
\end{tabularx}
\end{table}

\paragraph{What we found open}
Removal by expiry and exclusive leases both exist. Among the sources we read, none states
their combination for a hostile holder and disconnected verifiers, together with a lower
bound in the same model. This is weak evidence of a gap. We read five sources in full, and
we have not searched the forward citations of URSA and Nerio.
